\begin{ouverture}[geometrie]{14}{Trigonométrie du triangle rectangle}{Partie 4 · Espace et géométrie}
{Cosinus · Sinus · Tangente · Longueurs · Angles}
\objectif{Choisir le rapport adapté à l’angle et aux données.}
\objectif{Calculer une longueur ou un angle aigu.}
\objectif{Démontrer et utiliser $(\cos\widehat B)^2+(\sin\widehat B)^2=1$.}
\prerequis{Pythagore, Thalès et proportionnalité.}
\end{ouverture}
\pagestyle{otMain}\essentiel\fontsize{11}{15}\selectfont
\essSection{1}{Trois rapports pour un angle aigu}
\begin{definition}Dans ABC rectangle en A, pour l’angle $\widehat B$, BC est l’hypoténuse, AB le côté adjacent et AC le côté opposé. Les deux derniers rôles s’échangent si l’on choisit l’angle $\widehat C$.\end{definition}
\begin{propriete}$\cos\widehat B=\dfrac{AB}{BC}$, $\sin\widehat B=\dfrac{AC}{BC}$ et $\tan\widehat B=\dfrac{AC}{AB}$. Ces rapports ne dépendent que de l’angle : deux triangles rectangles ayant le même angle aigu peuvent être placés l’un dans l’autre, et le théorème de Thalès donne les mêmes rapports.\end{propriete}
\begin{savoirfaire}[Trouver une longueur]Écrire le rapport avec les noms des côtés avant de remplacer. Si BC = 10 cm et $\widehat B=35^\circ$, $AC=10\sin35^\circ\approx5{,}74$ cm. Si AB = 10 cm, on utiliserait $AC=10\tan35^\circ$.\end{savoirfaire}
\begin{propriete}Pour tout angle aigu, $(\cos\widehat B)^2+(\sin\widehat B)^2=1$. En effet, la somme vaut $\dfrac{AB^2+AC^2}{BC^2}$, et $AB^2+AC^2=BC^2$ d’après Pythagore.\end{propriete}
\begin{attention}La calculatrice doit être en mode degrés. Un sinus ou cosinus d’angle aigu est entre 0 et 1 ; une hypoténuse reste le plus long côté.\end{attention}
\clearpage
\essSection{2}{Angles et problèmes}
\begin{savoirfaire}[Trouver un angle]Si AB = 6 cm et BC = 10 cm, $\cos\widehat B=0{,}6$. Utiliser la fonction inverse du cosinus, souvent notée $\cos^{-1}$ ou acos : $\widehat B\approx53{,}1^\circ$. Cette touche ne calcule pas $1/\cos$.\end{savoirfaire}
\begin{propriete}Les deux angles aigus d’un triangle rectangle ont pour somme $90^\circ$. Cette relation fournit un contrôle ou un autre moyen de calcul.\end{propriete}
\begin{savoirfaire}[Une mesure indirecte]Pour une hauteur inaccessible, dessiner l’horizontale au niveau de l’observateur et sa ligne de visée. La hauteur calculée par tangente est celle au-dessus de cette horizontale ; ajouter ensuite la hauteur de l’instrument si nécessaire.\end{savoirfaire}
\begin{attention}Conserver les valeurs de calculatrice non arrondies jusqu’à la dernière étape. Préciser l’unité et la précision demandée. Un angle observé et une distance mesurés restent des données de modèle, avec une précision limitée.\end{attention}
\clearpage\rubrique{Découverte}
\begin{blocLivre}{ex-14-1}\begin{exercice}[1]{Nommer les côtés}
ABC est rectangle en A. Pour l’angle $\widehat B$, identifier le côté opposé, l’adjacent et l’hypoténuse.
\label{fig-ex-14-1}\figureGeo{trigo-01}
\end{exercice}\end{blocLivre}
\begin{blocLivre}{ex-14-2}\begin{exercice}[1]{Choisir le cosinus}
ABC est rectangle en A. AB = 6 cm et BC = 10 cm. Donner $\cos\widehat B$.
\label{fig-ex-14-2}\figureGeo{trigo-02}
\end{exercice}\end{blocLivre}
\begin{blocLivre}{ex-14-3}\begin{exercice}[1]{Choisir le sinus}
ABC est rectangle en A. AC = 8 cm et BC = 10 cm. Donner $\sin\widehat B$.
\label{fig-ex-14-3}\figureGeo{trigo-03}
\end{exercice}\end{blocLivre}
\begin{blocLivre}{ex-14-4}\begin{exercice}[1]{Choisir la tangente}
ABC est rectangle en A. AB = 6 cm et AC = 8 cm. Donner $\tan\widehat B$.
\label{fig-ex-14-4}\figureGeo{trigo-04}
\end{exercice}\end{blocLivre}
\begin{blocLivre}{ex-14-5}\begin{exercice}[1]{Une première longueur}
ABC est rectangle en A, $\widehat B=30^\circ$ et BC = 12 cm. Calculer AC.
\label{fig-ex-14-5}\figureGeo{trigo-05}
\end{exercice}\end{blocLivre}

\clearpage\rubrique{Entraînement}
\begin{blocLivre}{ex-14-6}\begin{exercice}[2]{Un côté adjacent}
ABC est rectangle en A. $\widehat B=40^\circ$ et BC = 9 cm. Calculer AB au millimètre.
\label{fig-ex-14-6}\figureGeo{trigo-06}
\end{exercice}\end{blocLivre}
\begin{blocLivre}{ex-14-7}\begin{exercice}[2]{Une hypoténuse}
ABC est rectangle en A. $\widehat B=50^\circ$ et AB = 7 cm. Calculer BC au dixième de centimètre.
\label{fig-ex-14-7}\figureGeo{trigo-07}
\end{exercice}\end{blocLivre}
\begin{blocLivre}{ex-14-8}\begin{exercice}[2]{Un angle par cosinus}
ABC est rectangle en A, AB = 6 cm et BC = 10 cm. Calculer $\widehat B$ au degré.
\label{fig-ex-14-8}\figureGeo{trigo-08}
\end{exercice}\end{blocLivre}
\begin{blocLivre}{ex-14-9}\begin{exercice}[2]{Un angle par tangente}
ABC est rectangle en A, AC = 5 cm et AB = 12 cm. Calculer $\widehat B$ au dixième de degré.
\label{fig-ex-14-9}\figureGeo{trigo-09}
\end{exercice}\end{blocLivre}
\begin{blocLivre}{ex-14-10}\begin{exercice}[2]{Cosinus et sinus liés}
ABC est rectangle en A. Démontrer que $(\cos\widehat B)^2+(\sin\widehat B)^2=1$. Sachant que $\cos\widehat B=0{,}6$, en déduire $\sin\widehat B$.
\label{fig-ex-14-10}\figureGeo{trigo-10}
\end{exercice}\end{blocLivre}

\clearpage\rubrique{Problèmes}
\begin{blocLivre}{ex-14-11}\begin{exercice}[3]{La hauteur d’un bâtiment}
Un appareil à 1,6 m du sol vise le sommet d’un bâtiment avec un angle d’élévation de 32 degrés. Il est à 25 m horizontalement du pied. Le sol est horizontal. Calculer la hauteur au décimètre.
\label{fig-ex-14-11}\figureGeo{trigo-batiment}
\end{exercice}\end{blocLivre}
\begin{blocLivre}{ex-14-12}\begin{exercice}[3]{Une rampe et son angle}
Une rampe monte de 0,8 m pour 10 m d’avancée horizontale. Dans ABC, A est l’angle droit, AB l’avancée et AC la hauteur. Calculer $\widehat B$ au dixième et la pente définie par $100\,AC/AB$.
\label{fig-ex-14-12}\figureGeo{trigo-12}
\end{exercice}\end{blocLivre}
\begin{blocLivre}{ex-14-13}\begin{exercice}[3]{Un câble}
Un câble BC de 18 m fait un angle $\widehat B$ de 65 degrés avec le sol horizontal AB. AC est vertical. Calculer la hauteur AC et la distance AB au dixième de mètre.
\label{fig-ex-14-13}\figureGeo{trigo-13}
\end{exercice}\end{blocLivre}
\begin{blocLivre}{ex-14-14}\begin{exercice}[3]{Un toit symétrique}
Le toit a pour section un triangle isocèle de base 8 m ; chaque versant fait 35 degrés avec l’horizontale. Calculer sa hauteur puis la longueur d’un versant au centimètre.
\label{fig-ex-14-14}\figureGeo{trigo-toit}
\end{exercice}\end{blocLivre}
\begin{blocLivre}{ex-14-15}\begin{exercice}[3]{Diagnostiquer une erreur}
ABC est rectangle en A. $\widehat B=38^\circ$, AB = 7 cm. Un élève annonce AC = $7\sin38^\circ$. Corriger puis donner AC au dixième.
\label{fig-ex-14-15}\figureGeo{trigo-15}
\end{exercice}\end{blocLivre}

\clearpage
